\subsection{Measures defined by numerical densities}

DEFINITION 2. --- A function f defined on T, with values in \(\overline{R}\) or in a Banach space, is said to be locally \(\mu\) -integrable if f is \(\mu\) -measurable and every point \(x \in T\) admits a neighborhood V such that \(\mu^{\bullet}(|\mathbf{f}|\varphi_{\mathrm{V}}) < +\infty\) .

This definition coincides with the one given in Ch. V, §5, No.~1, in case T is locally compact.

Let \(f\) be a locally \(\mu\) -integrable positive function; the mapping \(K \mapsto f_K \cdot \mu_K\) is a premeasure (Ch. V, §7, No.~1, Cor. 2 of Th. 1), which will be denoted \(f \cdot \mu\) .

PROPOSITION 2. --- If \(f\) is a positive locally \(\mu\) -integrable function, then, for every function \(g \in \mathcal{F}_{+}(\mathrm{T})\) , one has the relation

\[
(f \cdot \mu) ^ {\bullet} (g) = \mu^ {\bullet} (f g).\tag{1}
\]

Indeed, for every compact set K in T,

\[
(f \cdot \mu) _ {\mathrm{K}} ^ {\bullet} (g _ {\mathrm{K}}) = (f _ {\mathrm{K}} \cdot \mu_ {\mathrm{K}}) ^ {\bullet} (g _ {\mathrm{K}}) = \mu_ {\mathrm{K}} ^ {\bullet} (f _ {\mathrm{K}} g _ {\mathrm{K}}) = \mu_ {\mathrm{K}} ^ {\bullet} ((f g) _ {\mathrm{K}}),
\]

on using the definition of \(f \cdot \mu\) and Prop. 3 of Ch. V, §5, No.~3. Prop. 2 follows from this on passing to the supremum over \(K\) .

Now let \(x \in \mathbf{T}\) and let \(\mathbf{V}\) be a neighborhood of \(x\) such that \(\mu^{\bullet}(f\varphi_{\mathrm{V}}) < +\infty\) (Def. 2); then \((f \cdot \mu)^{\bullet}(\mathrm{V}) = \mu^{\bullet}(f\varphi_{\mathrm{V}}) < +\infty\) , therefore \(f \cdot \mu\) is a measure.

DEFINITION 3. --- Let \(f\) be a locally \(\mu\) -integrable positive function. The measure \(f \cdot \mu : K \mapsto f_K \cdot \mu_K\) is called the measure with density \(f\) with respect to \(\mu\) , or the product measure of \(\mu\) by the function f.~Every measure of the form \(f \cdot \mu\) , where f is positive and locally \(\mu\) -integrable, is called a measure with base \(\mu\) .

Remarks. --- 1) The definition of \(f \cdot \mu\) extends to the case that \(f\) is a complex locally integrable function; one then has \(|f \cdot \mu| = |f| \cdot \mu\) , which implies at once that \(f \cdot \mu\) is a measure, not just a premeasure. We retain the expression 'measures with base \(\mu'\) to designate the complex measures so defined.

\begin{enumerate}
\def\labelenumi{\arabic{enumi})}
\setcounter{enumi}{1}
\tightlist
\item
  Similarly, if \(\theta\) is a complex measure, \(f\) is said to be locally \(\theta\) -integrable if it is locally \(|\theta|\) -integrable, and one defines the measure \(f \cdot \theta : K \mapsto f_K \cdot \theta_K\) . One has \(|f \cdot \theta| = |f| \cdot |\theta|\) (Ch. V, §5, No.~2, Prop. 2). In this No., we shall leave aside everything concerning non-positive measures.
\end{enumerate}

PROPOSITION 3. --- Let \(\nu\) be a measure on T. For \(\nu\) to be of the form \(f \cdot \mu\) , where \(f\) is a locally \(\mu\) -integrable positive function, it is necessary and sufficient that every \(\mu\) -negligible compact set be \(\nu\) -negligible. If \(f'\) is a second locally \(\mu\) -integrable function such that \(\nu = f' \cdot \mu\) , then \(f = f'\) locally \(\mu\) -almost everywhere.

The condition is obviously necessary (Prop. 2). Conversely, suppose that every \(\mu\) -negligible compact set is \(\nu\) -negligible. Let us introduce a crushing \((\mathrm{K}_{\alpha})_{\alpha \in \mathrm{A}}\) of \(\mathbf{T}\) for the measure \(\mu + \nu\) and let us set \(\mathbf{N} = \mathbf{T} - \bigcup_{\alpha \in \mathrm{A}} \mathrm{K}_{\alpha}\) . It is clear that \((\mathrm{K}_{\alpha})_{\alpha \in \mathrm{A}}\) is a crushing for \(\mu\) and for \(\nu\) , and Prop. 9 of §1, No.~8 therefore implies the following relations for every \(g \in \mathcal{F}_+\) :

\[
\mu^ {\bullet} (g) = \sum_ {\alpha \in A} \mu_ {K _ {\alpha}} ^ {\bullet} \left(g _ {K _ {\alpha}}\right), \quad \nu^ {\bullet} (g) = \sum_ {\alpha \in A} \nu_ {K _ {\alpha}} ^ {\bullet} \left(g _ {K _ {\alpha}}\right).
\]

Consider a compact set \(C \subset K_{\alpha}\) that is \(\mu_{K_{\alpha}}\) -negligible; then C is locally \(\mu\) -negligible, hence locally \(\nu\) -negligible, and finally \(\nu_{K_{\alpha}}\) -negligible by the definition of \(\nu\) . It then follows from the Lebesgue--Nikodym theorem (Ch. V, §5, No.~5, Th. 2) that \(\nu_{K_{\alpha}}\) admits a density \(f_{\alpha}\) with respect to \(\mu_{K_{\alpha}}\) . Let f be the function that coincides with \(f_{\alpha}\) on each of the sets \(K_{\alpha}\) , and with 0 on N; the function f is \(\mu\) -measurable (Ch. IV, §5, No.~10, Prop. 16), and for every function \(g \in F_{+}\) one has, by the above relations and Prop. 3 of Ch. V, §5, No.~3,

\[
\nu^ {\bullet} (g) = \sum_ {\alpha \in \mathrm{A}} \nu_ {\mathrm{K} _ {\alpha}} ^ {\bullet} \left(g _ {\mathrm{K} _ {\alpha}}\right) = \sum_ {\alpha \in \mathrm{A}} \mu_ {\mathrm{K} _ {\alpha}} ^ {\bullet} \left(f _ {\alpha} g _ {\mathrm{K} _ {\alpha}}\right) = \sum_ {\alpha \in \mathrm{A}} \mu_ {\mathrm{K} _ {\alpha}} ^ {\bullet} ((f g) _ {\mathrm{K} _ {\alpha}}) = \mu^ {\bullet} (f g).
\]

It then follows first of all that \(f\) is locally \(\mu\) -integrable: if \(x\) is a point of \(T\) , and if \(V\) is a neighborhood of \(x\) such that \(\nu^{\bullet}(V) < +\infty\) , then \(\mu^{\bullet}(f\varphi_V) < +\infty\) . Next, Prop. 2 shows that the measures \(\nu\) and \(f \cdot \mu\) have the same essential upper integral. They are therefore equal (§1, No.~2, Cor. of Prop. 2). The uniqueness of \(f\) being obvious based on the case of compact spaces, the proposition is established.

Remark 3). --- The theory of integration with respect to a measure \(\nu = f \cdot \mu\) reduces at once to the theory treated in Ch. V. For, let \((\mathrm{K}_{\alpha})_{\alpha \in \mathbf{A}}\) be a crushing of T for \(\mu\) , hence for \(\nu\) , and let \(\mathrm{T}'\) be the locally compact space defined in the Scholium of §1, No.~8; we can associate to \(\mu\) (resp. to \(\nu\) ) a measure \(\mu'\) (resp. \(\nu'\) ) on \(\mathrm{T}'\) , in such a way that the measurable functions, the essentially integrable functions with values in a Banach space, and the essential upper integrals of positive functions, are the same for \(\mu\) and \(\mu'\) (resp. for \(\nu\) and \(\nu'\) ). The function \(f\) is therefore \(\mu'\) -measurable; it is locally \(\mu'\) -integrable, because \(\mathrm{T}'\) is locally compact, and a compact subset of \(\mathrm{T}'\) intersects only a finite number of the compact sets \(\mathrm{K}_{\alpha}\) ( \(\alpha \in \mathbf{A}\) ). Finally, the relation \(\nu'^{\bullet}(g) = \nu^{\bullet}(g) = \mu^{\bullet}(fg) = \mu'^{\bullet}(fg)\) proves that \(\nu' = f \cdot \mu'\) (Ch. V, §5, No.~3, Prop. 3). We leave to the reader the task of transcribing the results of Ch. V, §5.

